\section{A rational hitting list and its bit complexity}
\label{sec:generator}

Proposition~\ref{prop:pencil-deformation} supplies polynomial matrices
$R_i^*(t,p,r)$ for which $\sum_i A_i\otimes R_i^*$ is invertible
over $\Q(\omega)(t,p,r)$ whenever the pencil
$A_0+\sum_{i=1}^n A_i x_i$ is feasible and has order at most~$B$.
We now replace the root-of-unity coefficients and specialize the
three parameters to obtain a single list of rational matrix tuples.
The parameters $w,B,\Delta,M,N$ retain their earlier definitions.

To remove the root of unity, we represent multiplication in a
finite-dimensional rational algebra. Over $\C$ this representation
is a direct sum of root evaluations, so our invertibility conditions
must hold simultaneously at all roots. Set
\begin{equation}
  \ell=\frac M2,\qquad D=M\ell=\frac{M^2}{2},\qquad
  H=1+3(s+2)BM^3,
  \qquad
  \mathcal R=\Q[\xi]/(\xi^\ell+1).
  \label{eq:generator-parameters}
\end{equation}
Here and below $\xi$ also denotes its residue class. In particular,
$\xi^\ell=-1$ and $\xi^M=1$, so every integral power of~$\xi$, including
a negative power, reduces to a signed member of the basis
$1,\xi,\ldots,\xi^{\ell-1}$.

\subsection{The rational representation and the generator}

Let $\Omega$ be the set of complex roots of $X^\ell+1$. Since $M$ is a
power of two, these are the $\ell$ distinct numbers
\[
  \exp\!\left(\frac{2\pi\mathrm{i}(2j+1)}{M}\right),
  \qquad 0\le j<\ell.
\]
Every one is a primitive $M$th root of unity: its exponent $2j+1$ is
coprime to~$M$. Evaluation at these roots is an algebra isomorphism
\begin{equation}
  \mathcal R\otimes_{\Q}\C
    \longrightarrow \prod_{\omega\in\Omega}\C,
  \qquad f\longmapsto \bigl(f(\omega)\bigr)_{\omega\in\Omega}.
  \label{eq:cyclotomic-splitting}
\end{equation}
Indeed, a polynomial of degree less than $\ell$ is uniquely determined
by its values at $\ell$ distinct points, and interpolation gives a
polynomial with any prescribed values there.

For $a\in\mathcal R$, let $\rho(a)\in\Mat_\ell(\Q)$ be the matrix of
multiplication by~$a$ in the power basis. This is a unital algebra
homomorphism. It is computed using only rational arithmetic and the
relation $\xi^\ell=-1$. To make its splitting explicit, let $V$ be the
evaluation matrix whose $(\omega,j)$ entry is $\omega^j$, for
$\omega\in\Omega$ and $0\le j<\ell$. Interpolation says that $V$ is
invertible, and
\begin{equation}
  V\rho(a)V^{-1}
    =\operatorname{diag}\bigl(a(\omega):\omega\in\Omega\bigr).
  \label{eq:regular-representation-splitting}
\end{equation}
If $Q\in\Mat_d(\mathcal R)$, write $\rho_d(Q)$ for the rational
$d\ell$ by $d\ell$ matrix obtained by applying $\rho$ to every entry.
First changing basis by $I_d\otimes V$, and then reordering the basis
to group together coordinates belonging to the same root, gives
\begin{equation}
  \rho_d(Q)\sim_{\C}
    \bigoplus_{\omega\in\Omega}Q(\omega).
  \label{eq:matrix-regular-splitting}
\end{equation}
In particular, $\rho_d(Q)$ is invertible if and only if every
$Q(\omega)$ is invertible. These complex basis changes are used to
prove this assertion; the generator will not compute them.

We use the lift of Lemma~\ref{lem:division-lift}, now with $\xi$ in
place of a primitive root. For clarity, it can be computed directly
over~$\mathcal R$. Number the coordinates from $0$ to $M-1$ and set
\[
  D_\xi e_j=\xi^j e_j,\qquad Se_j=e_{j+1},
\]
with indices taken modulo~$M$. For $0\le i\le n$ define
\begin{align}
  \widetilde c_{i,a,b}(t)
    &=\frac1M\sum_{j=0}^{M-1}
          \xi^{-aj}(R_i(t))_{j,j-b},
       &&0\le a,b<M,\notag\\
  \widetilde R_i(t,p,r)
    &=\sum_{a,b=0}^{M-1}
          \widetilde c_{i,a,b}(t)p^a r^bD_\xi^aS^b
       \in\Mat_M\bigl(\mathcal R[t,p,r]\bigr).
  \label{eq:universal-cyclotomic-lift}
\end{align}
Evaluating its coefficients at $\xi=\omega$ gives exactly the lifted
matrix $R_i^{*,\omega}(t,p,r)$ from
Proposition~\ref{prop:pencil-deformation}. In particular, no choice of
a complex root is needed to compute~\eqref{eq:universal-cyclotomic-lift}.

The generator, on input $1^n,1^s$, performs the following finite
computation in a fixed order.
\begin{enumerate}
  \item Compute the integer parameters and the explicit matrices
    $R_0(t),\ldots,R_n(t)$ of
    Proposition~\ref{prop:rectangular}.
  \item For each $(t_0,p_0,r_0)\in\{1,\ldots,H\}^3$, compute
    \[
      Z_i=\rho_M\bigl(\widetilde R_i(t_0,p_0,r_0)\bigr)
      \in\Mat_D(\Q),\qquad 0\le i\le n.
    \]
  \item If $Z_0$ is singular, emit nothing for this triple. Otherwise
    emit the tuple
    \begin{equation}
      \bigl(Z_0^{-1}Z_1,\ldots,Z_0^{-1}Z_n\bigr).
      \label{eq:generator-output}
    \end{equation}
\end{enumerate}
Denote the resulting list by $\mathcal H_{n,s}$. Repetitions are
harmless. Every tuple has the same positive matrix dimension~$D$.

\subsection{The hitting property}

\begin{proposition}
\label{prop:generator-hitting}
Every nonzero admissible formula of size at most~$s$ has a defined,
invertible evaluation at a tuple in~$\mathcal H_{n,s}$.
\end{proposition}

\begin{proof}
Fix one such formula~$\Phi$. By
Proposition~\ref{prop:formula-certificates}, there is a family of at
most $s+1$ affine linear pencils over~$\Q$, each of order at most~$B$,
whose simultaneous invertibility ensures that the original formula
tree is defined and has an invertible value. These are the pencils
of all inverse-node occurrences and the root augmentation. Each has
an invertible evaluation over an extension field. Add the order-one
constant pencil~$1$ to this family, and call the resulting family
$\widehat{\mathcal C}_\Phi$. Thus
$|\widehat{\mathcal C}_\Phi|\le s+2$.

For $L=A_0+\sum_{i=1}^n A_i x_i\in\widehat{\mathcal C}_\Phi$ and
$\omega\in\Omega$, recall from Proposition~\ref{prop:pencil-deformation} that
\[
  P_{L,\omega}(t,p,r)
    =\det\!\left(\sum_{i=0}^n
          A_i\otimes R_i^{*,\omega}(t,p,r)\right).
\]
Proposition~\ref{prop:pencil-deformation} says that each of these is
a nonzero polynomial of total degree at most~$3BM^2$. Regard all of
them as polynomials in the one integral domain $\C[t,p,r]$ and form
\[
  P_\Phi(t,p,r)
    =\prod_{L\in\widehat{\mathcal C}_\Phi}
       \prod_{\omega\in\Omega} P_{L,\omega}(t,p,r).
\]
This is a finite product of nonzero polynomials, and
\begin{equation}
  \deg P_\Phi
    \le (s+2)\frac M2\,3BM^2
    =\frac32(s+2)BM^3
    <H.
  \label{eq:generator-product-degree}
\end{equation}
The extension field used to witness invertibility of a pencil is
not part of this coefficient ring: the $A_i$ are rational, and the
displayed lift has coefficients in $\Q(\omega)$. Thus there is no
need to embed that witness field into~$\C$.

A nonzero polynomial in three variables of total degree less
than~$H$ cannot vanish on all of $\{1,\ldots,H\}^3$. To see this,
suppose it did. Fixing the last two coordinates gives a univariate
polynomial of degree less than~$H$ with $H$ roots, so every
coefficient in the first variable vanishes on the two-dimensional
grid. Repeating the same argument for the second and then the third
variable makes every coefficient zero, a contradiction.
Choose a triple at which $P_\Phi$ is nonzero.

At this triple every root evaluation of $\widetilde R_0$ is
invertible, by the factor corresponding to the constant pencil~$1$.
Equation~\eqref{eq:matrix-regular-splitting} therefore shows that
$Z_0$ is invertible, so the generator retains the triple. The same
equation, applied to each raw pencil, shows that
\[
  \sum_{i=0}^n A_i\otimes Z_i
\]
is invertible for every $L\in\widehat{\mathcal C}_\Phi$.
If $L$ has order~$q$, left multiplication gives the exact identity
\begin{equation}
  (I_q\otimes Z_0^{-1})
       \left(\sum_{i=0}^n A_i\otimes Z_i\right)
   =A_0\otimes I_D+
       \sum_{i=1}^n A_i\otimes(Z_0^{-1}Z_i).
  \label{eq:generator-normalization}
\end{equation}
Hence all the certificate pencils are invertible on the rational
tuple~\eqref{eq:generator-output}.

The preceding splitting proved invertibility after extending
scalars to~$\C$. For a rational matrix this implies that its rational
determinant is nonzero; the adjugate formula then supplies its
inverse over~$\Q$. We may consequently apply
Proposition~\ref{prop:formula-certificates} directly over~$\Q$ to the
output tuple. It makes every required inverse of the original tree
defined over~$\Q$ and makes the final value invertible. This also
shows why no assertion about preservation of arbitrary inverse
domains under a homomorphism is needed.
\end{proof}

\begin{corollary}
\label{cor:pencil-hitting}
Let $L=A_0+\sum_{i=1}^n A_i x_i$ be a rational affine linear pencil
of order $q\le B$ that admits an invertible matrix evaluation over
an extension field of~$\Q$. There is a tuple
$X\in\mathcal H_{n,s}$ for which $L(X)$ is invertible.
\end{corollary}

\begin{proof}
Use the same argument with the two pencils $L$ and~$1$.
The product of their determinant polynomials over all
$\omega\in\Omega$ is nonzero and has total degree at most
$3BM^3<H$. A grid point where this product is nonzero is retained
by the generator, and~\eqref{eq:generator-normalization} makes
$L$ invertible at its output tuple.
\end{proof}

The proof of Proposition~\ref{prop:generator-hitting} fixes $\Phi$
only after the list has been specified. Its
polynomial $P_\Phi$ is used solely to prove the existence of a good
triple in the fixed grid; the generator never constructs this
polynomial, the certificates, or a witness evaluation. Different
formulas may use different triples. Arbitrarily large rational
constants change the coefficients of $P_\Phi$, but neither its
degree bound nor the construction of~$\mathcal H_{n,s}$.

\Needspace{10\baselineskip}
\subsection{Bit complexity and output length}

\begin{proposition}
\label{prop:generator-complexity}
The list $\mathcal H_{n,s}$ can be constructed and written by one
deterministic Turing machine in a number of bit operations bounded
by a fixed polynomial in $n+s+1$. Its entire binary encoding has
length bounded by a fixed polynomial in the same quantity.
\end{proposition}

\begin{proof}
Put $m=n+s+1$. The parameter definitions give
\[
  w\le m,\qquad B\le2m,\qquad
  \Delta\le4m^4,\qquad M\le2B\Delta\le16m^5.
\]
Consequently
\begin{equation}
  H=O(m^{17}),\qquad D=O(m^{10}),\qquad
  H^3=O(m^{51}).
  \label{eq:generator-dimension-bounds}
\end{equation}
Counting the unary inputs, forming these parameters, and finding
$M$ by successive doubling all take polynomial time.

We first bound the integers actually computed at a grid triple.
For an input monomial $z^j$, where $0\le j<N$, the defining formula
for $T_i$ gives
\begin{equation}
  T_i z^j
    =\sum_{a=0}^j \binom ja i^{j-a}t^a(z-i)^{M-1-a}.
  \label{eq:generator-binomial-formula}
\end{equation}
Expanding each power in the monomial basis computes every entry
using integer arithmetic. Each coefficient in either of the two
basis expansions has absolute value at most $2^M w^M$. At a grid
value $1\le t\le H$, an entry of the padded matrix $R_i(t)$
therefore has absolute value at most
\[
  M\,2^{2M}w^{2M}H^M.
\]
In~\eqref{eq:universal-cyclotomic-lift}, the only division of a
coefficient is by~$M$. All powers of $\xi$ reduce to signed basis
monomials. The Fourier sum and the lift involve polynomially many
terms, and $p^a r^b\le H^{2M}$ on the grid. For example, the
absolute values of the power-basis coefficients of every entry of
$M\widetilde R_i(t,p,r)$ are bounded by
\[
  M^4\,2^{2M}w^{2M}H^{3M}.
\]
Multiplication by a power-basis monomial merely permutes those
coefficients and changes signs. Thus
\begin{equation}
  U_i:=MZ_i\in\Mat_D(\mathbb Z),\qquad
  b=O\!\left(M\log\bigl((w+1)(H+1)\bigr)+\log M\right)
    =O(m^5\log m)
  \label{eq:generator-input-bit-bound}
\end{equation}
is a common bound on the bit lengths of their entries. These
calculations have explicit nested loops of polynomial length:
binomial coefficients may be formed by Pascal's recurrence,
integer powers by repeated multiplication, and all ring
coefficients by the signed reductions just described. Every
intermediate sum and product in these loops has polynomial bit
length. The matrices $U_i$ can therefore be constructed in
polynomial bit time.

The common clearing factor cancels in normalization:
\begin{equation}
  Z_0^{-1}Z_i=U_0^{-1}U_i.
  \label{eq:generator-common-denominator}
\end{equation}
If $U_0$ is nonsingular, Cramer's rule computes each column of this
matrix using $\det U_0$ and determinants of $U_0$ with one column
replaced by a column of~$U_i$. All these are integer matrices of
order~$D$ with the same entry bit bound~$b$. Every minor of any such
matrix has absolute value at most $D!2^{Db}$. Hence its bit length
is bounded by
\begin{equation}
  \Lambda=O\bigl(D(b+\log(D+1))\bigr)
     =O(m^{15}\log m).
  \label{eq:generator-minor-bit-bound}
\end{equation}
Each output entry is a ratio of two such determinants. It can be
written with a signed binary numerator and a positive binary
denominator using $O(\Lambda)$ bits, changing both signs if
necessary. Cancelling their gcd is optional for this encoding.

For completeness, exact Gaussian elimination computes the needed
determinants in polynomial bit time without permitting fractions
to grow unchecked. Let $A$ be any one of these integer matrices.
After $k$ successful pivots, let $I_k$ be the ordered list of the
selected original row indices, and let
\[
  B_k=A[I_k,\{1,\ldots,k\}].
\]
This submatrix is invertible. For an unselected row index $u$ and
a column $v>k$, the remaining Gaussian entry is
\begin{equation}
  A_{uv}-A[u,\{1,\ldots,k\}]\,
         B_k^{-1}A[I_k,v]
   =\frac{
       \det\begin{pmatrix}
         B_k&A[I_k,v]\\
         A[u,\{1,\ldots,k\}]&A_{uv}
       \end{pmatrix}}
       {\det B_k}.
  \label{eq:generator-pivot-minors}
\end{equation}
For $k=0$ the denominator is~$1$ and the entry is $A_{uv}$.
The determinant identity follows by subtracting the appropriate
linear combination of the first $k$ block rows from the last row.
It also shows inductively that ordinary Gaussian updates compute
the displayed entries.

If the next pivot column has a nonzero remaining entry, swap its
row into the pivot position and append that original row index to
$I_k$. These swaps involve only unselected rows, so they do not
change any previously selected submatrix. Writing the next pivot
as~$p_{k+1}$, equation~\eqref{eq:generator-pivot-minors} gives
\[
  p_{k+1}=\frac{\det B_{k+1}}{\det B_k},
  \qquad p_1\cdots p_k=\det B_k.
\]
If the entire next remaining column is zero, the Schur complement
is singular and $\det A=0$. Otherwise all $D$ pivots can be taken,
and their product, multiplied by the sign of the accumulated row
permutation, equals~$\det A$.

Both the stored Gaussian entries and the accumulated pivot
products, when reduced to fractions, therefore have
$O(\Lambda)$-bit numerators and denominators: they are ratios of
minors or minors themselves. Each individual update combines only
a constant number of these fractions, so its intermediate products
and sums also have $O(\Lambda)$ bits before reduction. Integer
addition, multiplication, and division on such operands have
polynomial bit cost by the elementary digit algorithms. Gcd
reduction does as well: in the Euclidean algorithm the nonzero
remainders decrease by a factor of at least two within every two
steps, so there are $O(\Lambda)$ polynomial-cost divisions.
There are $O(D^3)$ Gaussian arithmetic updates. Pivot searches,
row swaps, and sign bookkeeping also have polynomial cost.
This proves polynomial bit time for each determinant, including
singular matrices and all the column-replaced matrices used by
Cramer's rule.

At each grid triple, at most $nD^2+1$ determinants suffice for the
singularity test and all output entries. Equations
\eqref{eq:generator-dimension-bounds}--\eqref{eq:generator-minor-bit-bound}
therefore give a fixed polynomial bound on the total computation.
The complete list contains at most $H^3 nD^2$ rational entries, so
its encoding length is
\begin{equation}
  O(H^3 nD^2\Lambda)=O(m^{87}\log m).
  \label{eq:generator-output-bit-bound}
\end{equation}
Dimensions, signs, and fixed delimiters fit within the same bound.
A fixed delimited encoding allows the list to be written as its
entries are produced; if its length is placed first, one may instead
buffer the polynomial-length list or make a preliminary counting
pass.

Finally, these arithmetic and loop bounds give a Turing-machine
bound as well. All arrays have polynomially many entries of
polynomial bit length. They may be stored as indexed binary records
on work tapes; even scanning all records for each requested access
adds only a fixed polynomial factor. The loops, elementary digit
algorithms, and delimiters are fixed in advance. Thus one
deterministic Turing machine performs the construction and writes
the whole list within $C(n+s+1)^k$ bit operations for some fixed
positive integers $C,k$. None of its arithmetic uses the constants
or coefficients of an unknown formula.
\end{proof}

\begin{proof}[Completion of the proof of Theorem~\ref{thm:main}]
Use the generator just defined. Proposition~\ref{prop:generator-hitting}
gives the universal hitting property, with invertible and therefore
nonzero output values. Proposition~\ref{prop:generator-complexity}
gives both required polynomial bounds, after increasing the same
fixed constants $C,k$ if necessary. The generator is deterministic
and depends only on its two unary inputs.
\end{proof}
